\cleardoublepage
\newrefsection
\chapter{文献综述}

\section{背景介绍}
大语言模型（Large Language Models, LLMs）\cite{transformer, gpt3, llama}的迅猛发展彻底改变了自然语言处理（NLP）领域的格局，GPT-4\cite{gpt4}、Claude 3\cite{claude3}等模型在多项任务中展现出了接近人类的理解与生成能力。然而，LLM仍面临着诸如“知识截止”（Knowledge Cutoff）、“幻觉”（Hallucination）\cite{ref24}、领域专业知识匮乏以及无法访问私有数据等关键挑战。检索增强生成（Retrieval-Augmented Generation, RAG）\cite{ref1}作为一种结合了预训练参数化记忆（Parametric Memory）与外部非参数化记忆（Non-Parametric Memory）的范式，通过在生成过程中动态检索并利用外部知识库，有效地缓解了上述问题，显著提升了生成内容的准确性与时效性。

尽管RAG技术在简单问答任务中表现优异，但在面对科学文献综述、医学诊断、法律案件分析等复杂场景时，传统的单跳检索（Single-hop Retrieval）往往难以奏效。这类任务通常涉及跨文档的知识整合与多步推理（Multi-hop Reasoning），需要模型具备在海量异构数据中追踪线索、关联事实的能力。此外，现实世界中的知识并非单一的纯文本形式，而是广泛分布于非结构化文本、半结构化表格、结构化知识图谱（Knowledge Graph, KG）以及图像等多模态数据中。如何打破数据孤岛，实现多源异构知识的深度融合与高效推理，已成为当前人工智能与知识工程交叉领域的关键科学问题与研究热点。

从系统工程角度看，RAG落地过程中还存在两类常被忽视但对可用性影响显著的问题：其一，检索到的证据并不总是“可用证据”，噪声候选、过长上下文与证据碎片化会导致生成端对关键信息的利用效率下降；其二，真实用户查询包含大量歧义、否定约束与无答案情形，若系统缺少对证据充分性的判断与拒答机制，则容易产生高置信的错误回答，从而削弱用户信任。因此，当前研究不仅需要提升平均指标，更需要关注在困难输入分布下的鲁棒性与可解释性，建立从检索质量、排序稳定性到事实核验的全链路评测与迭代方法。

\section{国内外研究现状}

围绕“自适应检索”“多源知识融合”“复杂推理”三条主线，本文重点检索并归纳了近三年（2023--2025）发表在CCF A类会议/期刊及arXiv的前沿工作，并按技术范式进行分类对比。为保证综述具有代表性与可比性，本文选取并重点对比了10项核心工作：Self-RAG\cite{ref8}、Corrective RAG\cite{ref9}、Active RAG\cite{ref12}、Adaptive RAG\cite{ref11}、RankRAG\cite{ref10}、RAPTOR\cite{ref13}、GraphRAG\cite{ref15}、LightRAG\cite{ref16}、KG-Agent\cite{ref18}、G-Retriever\cite{ref33}，并在此基础上结合经典方法与评测指标进行补充讨论。

\subsection{RAG基础架构演进}
检索增强生成技术自提出以来，经历了从“检索—拼接—生成”的基础形态逐步演进到更强调控制与可扩展性的系统化形态。早期的Naive RAG通常将检索到的若干文档片段直接拼接进提示词并交由模型生成，这一流程在一定程度上扩展了模型的知识边界，但其效果高度依赖检索器的召回质量：一旦检索结果噪声较高或遗漏关键证据，生成质量会明显下降；同时受限于上下文窗口长度，该类方法难以在长文档或高召回量场景下维持稳定表现。

为缓解上述问题，Advanced RAG在“检索前—检索后—输入构造”三个环节引入了更细粒度的优化。检索前侧重通过查询重写（Query Rewriting）与查询扩展（Query Expansion）减少歧义并补全意图；检索后侧重利用重排序（Re-ranking）模型对候选文档进行精排，结合上下文压缩（Context Compression）或信息抽取将有限窗口尽可能分配给高密度、与问题强相关的证据，从而提升证据利用效率与答案一致性\cite{manning_ir, rocchio, bert_rerank, llmlingua}。

随着RAG系统逐渐承担更复杂的任务形态，模块化趋势开始显现。Modular RAG\cite{ref27}将RAG拆分为检索、生成、评估、路由等相对独立的组件，并通过标准化接口支持不同模块的组合与替换，使系统能够针对领域数据与任务需求进行定制化优化，同时也为后续引入自适应检索、复杂推理与可信验证等能力提供了工程载体。

在架构演进的过程中，“排序与选择”逐渐成为影响RAG效果的关键控制点。一方面，密集检索器（如DPR）\cite{ref6}在语义召回上具备优势，但往往需要依赖重排序器对Top-$k$候选进行精排，才能稳定提供高质量证据；另一方面，随着生成模型上下文窗口增长，简单拼接并不能保证证据被有效利用，如何在有限预算下进行证据选择、证据聚合与证据压缩，成为系统性优化的重要方向。相关工作也开始把排序信号融入端到端框架，强调“检索—排序—生成”作为一个整体需要联合评估与调参\cite{ref10}。

\subsection{自适应检索机制}
传统RAG往往对每个查询都执行一次固定检索，但并非所有问题都需要外部证据支持，过度检索既增加计算开销，也可能引入噪声干扰推理。因此，近年来研究开始把“何时检索、检索多少、检索哪些来源”作为可学习或可控制的决策问题来处理，自适应检索（Adaptive Retrieval）由此成为重要方向。

从技术范式上看，近三年的自适应检索工作可归纳为三类：\textbf{（1）反思式/自评式控制}，代表如Self-RAG\cite{ref8}在生成过程中插入反思信号，对“是否需要检索、检索证据是否相关”进行自评，从而实现按需检索与证据过滤；\textbf{（2）质量感知的纠错闭环}，代表如Corrective RAG\cite{ref9}显式评估检索证据质量，当证据不足或偏离时触发纠错检索或外部搜索，使系统在困难分布上更鲁棒；\textbf{（3）策略路由与主动检索}，代表如Adaptive RAG\cite{ref11}以查询复杂度/不确定性为信号进行路由，Active RAG\cite{ref12}把检索触发与生成过程交错，使检索从“固定前置步骤”变为“过程内行动”。此外，RankRAG\cite{ref10}将排序信号显式纳入RAG框架，为“检索后怎么用证据”提供了统一接口；RAPTOR\cite{ref13}以树组织文档摘要与检索，体现了在长文档与多粒度证据下的结构化检索思路。

总体而言，自适应检索的优势在于能够以更少的无效检索换取更高的证据命中率，并在证据不足时显式触发纠错或扩展检索，从而提升困难样本上的鲁棒性\cite{ref8, ref9, ref11, ref12}；其局限也较为明确：检索决策往往依赖模型自评或启发式信号，稳定性与可迁移性仍不足；同时，多轮检索与重排序会显著增加端到端延迟与调用成本，若缺少对“调用次数、时间与Token开销”的约束与评测，自适应检索容易停留在“效果提升但不可用”的状态\cite{ref10, ref12}。

自适应检索的进一步挑战在于：检索决策不仅要考虑是否检索，还要考虑“检索后怎么用”。当检索结果本身不稳定或证据分布高度分散时，系统需要决定是否触发二次检索、是否扩大召回范围、是否引入不同检索范式（稀疏/稠密/图遍历）以及如何在多个候选证据之间进行加权融合。换言之，自适应检索更像是一个贯穿全流程的控制问题，其信号来源可以是模型不确定性、检索分数分布、证据一致性甚至用户反馈。围绕这些信号构建可解释的策略与可复现的评测，将直接决定自适应检索是否能够从“经验技巧”走向“工程可用”。

\subsection{多源异构知识融合}
在更贴近真实世界的应用中，知识往往分布于多种载体：非结构化文本提供丰富语义与叙事信息，表格与数据库承载结构化事实与数值关系，知识图谱刻画实体关系与可遍历逻辑。单一模态检索很难同时覆盖“语义匹配、结构约束与关系推理”的需求，因此多源异构知识融合成为提升RAG深度与可解释性的关键路径。

围绕“文本—图谱”的融合，GraphRAG\cite{ref15}通过LLM从文本中抽取实体与关系并构建图结构，再以社区发现等方法形成不同层级的摘要与索引，适合回答需要全局结构与宏观概括的问题。LightRAG\cite{ref16}进一步强调效率，通过双层检索与图结构索引在一定精度下提升处理速度；HippoRAG\cite{ref14}借鉴海马体记忆机制，以个性化PageRank\cite{pagerank}等方式在图上模拟联想检索，支持跨段落的隐式关联发现。与此同时，Think-on-Graph\cite{ref17}让模型在图上执行游走与推理，减少结构化知识线性化带来的逻辑断裂；KG-Agent\cite{ref18}则以工具调用方式把自然语言查询拆解为结构化查询语句，从而对大规模知识库进行更精准的访问与验证。

除图谱外，表格等载体同样是重要的知识来源。Hybrid RAG\cite{ref19}探索了语义检索与关键词检索的混合策略，并强调在混合检索结果上进行统一排序与融合。面向“文本+表格”的问答任务，HybridQA\cite{hybridqa}、OTT-QA\cite{ottqa}等基准推动了跨载体证据组合的研究，也提示系统需要同时处理自然语言语义匹配与表格字段/行列约束。总体来看，异构融合研究正在从“把不同模态拼接在一起”走向“分工检索 + 证据对齐 + 互证约束”的系统化方向：先在不同载体内以最合适的表示与检索器获得候选证据，再在统一的证据空间中完成对齐、冲突消解与融合。

这一路线的优势在于能够显式利用图谱路径与表格字段约束，提供更强的可解释结构，并在全局层面组织证据以支撑多跳推理\cite{ref15, ref16, ref17, ref18}；但其局限也同样突出：图构建与实体对齐依赖抽取质量与消歧能力，误差会在后续推理中被放大；表格与文本之间存在天然的结构差异，序列化与向量化往往会损失行列逻辑与数值比较信息，导致“可检索但不可推理”；此外，多源证据带来的冲突与版本不一致问题仍缺少统一、可复现的对齐与核验范式\cite{hybridqa, ottqa, ref19}。

值得注意的是，异构融合并不等同于简单的“多路召回”。当系统同时引入文本、表格与图谱时，不同证据源的对齐与一致性将成为核心难点：例如同一实体在不同载体中的别名、时间版本差异、表格字段语义与文本叙述的不一致等，都可能导致推理出现偏差。因此，越来越多的工作开始关注“统一表示 + 结构约束”的组合，即在共享向量空间中完成粗筛，再利用结构化关系、字段约束或路径一致性进行精筛与验证。这种思路也为后续的可解释性提供了基础，因为证据不仅可以以文本片段呈现，还可以以结构路径与字段匹配的形式被审计。

\subsection{复杂推理机制}
推理能力是面向复杂任务的核心能力之一。为增强推理，研究首先从“让模型显式写出中间步骤”出发，Chain-of-Thought（CoT）\cite{ref2}通过提示引导模型生成推理链条，提升逻辑一致性与可控性；随后，Tree of Thoughts（ToT）\cite{ref4}引入树搜索思想，使模型能够探索多条推理路径并在必要时回溯修正；Graph of Thoughts（GoT）\cite{ref5}进一步将推理组织为更一般的图结构，以支持并行探索、信息聚合与循环反馈，从而更好地处理多因素依赖与复杂约束。

在RAG系统中，推理与检索的关系也从“检索先行、推理随后”逐渐转向深度协同。IRCoT\cite{ref26}提出在推理过程中交错地执行检索步骤，即“检索—推理—再检索—再推理”，使每一步推理都能获得最新证据支持，从而缓解多跳问答中的信息断层。与此同时，ReAct\cite{ref23}将推理与行动结合，允许模型在推理过程中调用外部工具（搜索、计算、API等），为复杂任务提供可执行的证据获取与验证途径，也为构建可解释、可审计的系统行为带来新的工程形态。

从结果可靠性的角度看，复杂推理不仅要求“推理能走通”，还要求“推理可被证据约束与检验”。当模型输出推理步骤时，如何确保每一步都能在检索证据中找到支撑，如何在证据缺失时及时回退或拒答，是复杂推理落地的关键。围绕这一目标，一类思路是在生成后进行链式核验\cite{ref24}，另一类思路是对生成内容进行细粒度事实评分\cite{ref25}，并将核验结果反哺到检索与推理控制中。把推理结构、检索证据与核验机制统一到同一条评测链路，将有助于形成可持续迭代的“可信RAG”系统。

从方法对比看，结构化思维范式（CoT/ToT/GoT）能够显式暴露中间步骤并支持回溯，但其“推理文本”未必对应外部证据，容易出现看似合理的推理幻觉；交错检索范式（如IRCoT）能把推理与证据获取绑定，提升多跳任务的可达性，但会引入更多检索轮次与更高延迟；工具调用与核验范式（如ReAct、CoVe与事实评分）能把系统推向可审计与可控，但需要设计稳定的触发信号与评测口径，避免过度核验导致成本不可控\cite{ref2, ref4, ref5, ref23, ref24, ref25, ref26}。

\section{研究展望}

尽管现有研究在检索、重排序、图谱融合与推理结构等方面取得了显著进展，但面向真实复杂场景的统一框架仍处于探索阶段。结合近三年前沿工作对比，当前研究可以进一步凝练为三个关键问题：\textbf{（1）异构证据的可对齐融合问题}，系统既要覆盖文本、表格与图谱等多源知识，又要保证同一事实在不同载体间可定位、可互证与可审计；\textbf{（2）自适应检索的成本与可控性问题}，按需检索往往带来多轮调用与额外推理开销，如何在保证效果的同时降低延迟、控制调用次数并提升策略稳定性，是工程落地的核心约束；\textbf{（3）复杂推理的可验证闭环问题}，推理过程需要被证据链约束并支持核验与回退，避免在证据不足时产生高置信错误输出，从而形成可持续迭代的可信RAG评测闭环。

未来的研究应致力于构建面向多源异构知识的RAG复杂推理框架，探索异构知识的统一表示与索引，设计自适应多策略检索机制，并增强推理过程的可解释性与验证机制，从而推动人工智能系统在复杂真实场景中的落地应用。

针对上述问题，本文的研究思路可以概括为以“证据”为中心的系统化设计：在数据与索引层面实现异构载体的统一入口，在检索与排序层面引入可解释的控制策略，在推理层面以图结构组织中间状态并保持可回溯，在验证层面通过核验与拒答降低高置信错误输出。通过建立统一的评测闭环与迭代机制，本研究希望为面向真实复杂场景的RAG系统提供一条兼顾效果、成本与可信性的可落地路径。

\newpage
\begingroup
    \linespreadsingle{}
    \printbibliography[title={参考文献}]
\endgroup
